\section{2026 Active Aero Investigation}
\label{sec: 2026 Active Aero Investigation}

The 2026 F1 regulations\cite{2026_FIA_regs} replace DRS with an automatic active aero system operating without any gap-to-car-ahead requirement. This removes the slipstream confound and, in principle, enables a clean $\Delta C_DA$ measurement from the drag difference between straight mode (active aero open) and corner mode (active aero closed) coast-down segments within the same free-practice session.

Canada 2026 Race and FP1 telemetry (all 22 drivers) were retrieved via FastF1 and inspected for active aero state. Channel 45 (the DRS indicator in 2024, values 0/12/14) is absent from every 2026 car data message. The FastF1 source code acknowledges this: \texttt{entry['Cars'][drv]['Channels'].get('45',\,0)} with a comment \textit{`drs is no longer included in 2026'}. No replacement channel carrying active aero state is present in the 2026 API.

Without active aero state telemetry, the natural fallback is to classify segments by track position using published active aero zone boundaries. This is a methodological challenge, not a fundamental blocker: with precisely known zone boundaries, a conservative exclusion buffer around the transition points, and restriction to segments lying entirely within a single aerodynamic mode, position-based classification could yield usable coast-down pairs. The drag difference between a straight-mode segment (active aero open, high speed) and a corner-mode segment (active aero closed, lower speed, same session) would then be attributable to $\Delta C_DA$, free of the slipstream confound that made the DRS analysis inconclusive.

\textbf{Proposed position-buffer approach.} A buffer of $\pm 75$~m around each published zone boundary (assuming zone boundaries are published in the FIA Event Notes, as was the case for 2024 DRS zones; if these documents do not cover 2026 active aero zones, this approach requires a verified boundary source before implementation) would be excluded, leaving segments that lie entirely within a confirmed mode. For Canada Circuit Gilles Villeneuve, the two DRS zones (pit straight and back straight) each span $\sim\!650$~m, so a $\pm 75$~m buffer removes roughly 23\% of each zone---still leaving $\sim\!500$~m of usable straight. Sensitivity to buffer width (25--150~m) would quantify how much $\hat{\alpha}$ changes as transition-contaminated samples are progressively excluded; convergence of $\hat{\alpha}$ across buffer widths would provide indirect evidence that the chosen buffer is adequate.

The genuine hard blocker is the unknown command timing. The active aero system transitions between modes as the car crosses a zone boundary, but the exact moment of actuation can vary by several tenths of a second depending on vehicle speed, system latency, and team calibration. At 80~m/s (288~km/h) a 0.3~s timing uncertainty corresponds to 24~m of position error, comparable to a $\pm 25$~m buffer---making the smallest buffers unreliable. Without knowing when the transition occurs, the exclusion buffer cannot be sized with confidence, and transition-contaminated samples may be included in the wrong mode class.

The measurement will become straightforwardly feasible once the F1 timing API exposes the active aero state channel. A position-only fallback using a $\geq\!75$~m buffer with sensitivity testing on buffer width is a viable interim approach and is recommended as the next implementation step.